\chapter[Conception and Improvement]{Basics for Conception and Improvement}
\label{sec:BasicImprovement}
\newpage
\section{Basic shapes}	
	\label{sec:BasicShapes}
\subsection*{The ``traditional'' shape}
\begin{center}
	\includegraphics[width=30em]{22_BRECH.eps} 
\end{center}
The ``traditional'' shape is the simplest. With this you can achieve anything from long distance flights to short range low, rapid flights, to long flight durations. Since both wing tips are relatively far away from the centre of rotation, however, these designs tend to maintain their ``spinning momentum'', making it easier to injure your fingers when catching.
\newpage

\subsection*{The hook shape}
These are more or less the basic forms from which all other boomerang shapes are derived. 

\begin{center}
	\includegraphics[width=30em]{23_HOOK.eps} 
\end{center}
The hook form is especially designed for long distance. It is more difficult to catch, because of the wide inner-circle yet less painful than the standard form.
\newpage

\subsection*{The multiblade shape}
\begin{center}
	\includegraphics[width=30em]{24_DFL.eps} 
\end{center}
The multiblader shapes tend to have shorter range, but retain their rotation and stability to the end of the flight. Catches are easier since there is material at the centre of rotation.
\newpage

\section{Airfoils}
Choice of airfoil depend on the flight characteristics desired. The following airfoils are considered as basic airfoils:
\begin{center}
\includegraphics[width=18em]{25_PSTD.eps}\\
Standard airfoil\\
\vspace{2em}
\includegraphics[width=18em]{26_PRODU.eps} \\
Completely profiled airfoil\\
\vspace{2em}
\includegraphics[width=18em]{27_STDBE.eps}\\
Standard airfoil with undercut\\
\vspace{2em}
\includegraphics[width=18em]{28_PRAUS.eps} \\
Hollowed out standard airfoil\\
\vspace{2em}
\includegraphics[width=30em]{29_PRON.eps} \\
``Neutral'' airfoils e.g. on the elbow\\
\vspace{2em}
\includegraphics[width=18em]{30_SYMLR.eps}\\
Airfoil for right- and left-handers\\
\end{center}               
\newpage

\begin{center}
\vspace{2em}
\includegraphics[width=18em]{31_SYMOU.eps} \\
Reverse airfoil (the boomerang can be thrown upside-down)\\
\vspace{2em}
\includegraphics[width=30em]{32_SCHLI.eps} \\
Slit airfoil
\end{center}

The following airfoils are merely recommendations for various types of boomerang.
\begin{center}
\includegraphics[width=30em]{34_MTA.eps} \\
MTA - airfoils\\
\vspace{2em}
\includegraphics[width=18em]{35_PROLD.eps} \\
Long Distance airfoil\\
\vspace{2em}
\includegraphics[width=18em]{36_PFC.eps} \\
Fast Catch airfoil\\
\vspace{2em}
\includegraphics[width=18em]{38_WIND.eps} \\
Strong wind airfoil\\
\vspace{2em}
\end{center}
\newpage

\section{Material}
The best material, as already mentioned, is finnish birch plywood as it is easy to work with, tough and flexible. Boomerangs can also be made out of other plywoods, solid and balsa wood, as well as paper, cardboard and various synthetics such as plexiglass, polypropylene, ABS, phenolic and fibreglass or carbon fibre reinforced plastic. The toughness and density of the material is also of interest. Choose the appropriate material for the boomerang design you are making. Costs of materials vary considerably. All the materials have one thing in common: they are all strong yet flexible and fairly tough. Finnish birch plywood is still the best compromise for beginners and for ``normal'' boomerangs, even though polypropylene and ABS is becoming more and more popular.

For example, for MTA boomerangs, the material should be between 1.5 mm and 3 mm thick. If you use wood it should be 3 mm thick. Polypropylene, 3 mm thick is not stiff enough. carbon fibre, fibreglass or plexiglass are too heavy. 1.5 - 2.0 mm Phenolic is almost ideal. It is strong enough, without being too heavy. It is fairly tough but a touch too heavy. 
For Long Distance boomerangs, carbon fibre, fibre glass, phenolic and wood work well. Fibreglass and carbon fibre tend to be more expensive.
\newpage

\begin{center}
\colorbox{BOrange!30}{\parbox{30em}{\centering
\begin{center}
\textbf{Facts on Materials}\\
\end{center}
\colorbox{BOrange!30}{\parbox{28em}{
	\color{black}
	\vspace{-0.5em}
Please be aware of the safety instructions in section \ref{Sec:SafetyMaterials}!
\paragraph*{Finnish birch plywood:} 2 layers of wood per mm, should be at least 3 mm thick, easily bendable, compact/density: approx. 0,7 Kg/dm3
\paragraph*{Phenolic:} paper dipped in phenol resin, should be at least 2 mm thick, soft, easily bendable, density approx. 1,5 Kg/dm3
\paragraph*{PP:} polypropylene, at least 3 mm thick, very soft, easily bendable, unbreakable, density: approx. 0,92 Kg/dm3
\paragraph*{ABS:} durable polystyrol mixed with butadien and acylic nitril, at least 3 mm thick, hard, easily bendable, relatively fragile, density approx. 1,05 Kg/dm3
\paragraph*{Plexiglass:} acrylic glass or polymethylmetha-crylat, usable at a thickness of 3 mm, soft, difficult to bend, unbreakable, density: 1,18 Kg/dm3
\paragraph*{Fibre glass:} Glass fibre synthetic, usable at a thickness of 1,5 mm, hard, bendable, fairly tough, density: 1,8 Kg/dm3
\paragraph*{Carbon fibre:} carbon fibre synthetic, usable at a thickness of 1,5 mm, hard, bendable, fairly tough, density; approx. 1,8 Kg/dm3
	\vspace{0.5em}
}}}}
\end{center}

\section{Surface}
The texture of the surface can influence the boomerang’s flight pattern. A rough surface tends to reduce hover at the end of the flight. A boomerang with a smooth surface flies faster, retains its spin, and tends to hover more.
More lift can be produced by a smooth topside and a rough underside. For reverse circumstances (rough topside and smooth underside) the lift is decreased. The increase in lift is, however, not as significant as an increase in positive angle of attack, for example.
\newpage 

\section{Tuning (Bending and twisting)}
There are two methods of bending a boomerang. The first is bending a wingtip up to increase ``dihedral'' (see figure \ref{fig:PosDihedral}). The second method is twisting the leading edge of the wing upward to increase angle of attack (see figure \ref{fig:AngleAttack}). 
\subsection*{Dihedral}
\label{sec:dihedral}
Generally speaking, a slight positive dihedral in both arms, stabilizes the flight of a boomerang. If there is too much dihedral the boomerang rises too high and tends to have an S-shape flight path. A boomerang with a negative dihedral flies further and lower. If it has been tuned too negative, the boomerang is most likely to fly straight ahead and then crash to the ground. 

\begin{figure}[H]
\centering
\includegraphics[width=30em]{39_LBIEG.eps} 
\caption{Positive Dihedral}
\label{fig:PosDihedral}
\end{figure}

With the help of the construction disk (see p. 20), you can see what happens (wing tips gather more forward speed ($V_t$) than rotation speed ($V_r$)). The positive dihedral creates maximum lift at Points $B_1$ and $B_2$. The resulting turning force occurs $90^\circ$ later at $B_1^{\prime}$ and $B_2^{\prime}$.

When arm 1 has positive dihedral, the boomerang has more lay-over and will fly a bit further. Positive dihedral on arm 2 just tends to cause more lay-over. Positive dihedral on arm 2 has a larger effect than arm 1 because it flies through the air at a higher speed (see figure \ref{fig:DiVtVr}). The higher the forward speed ($V_t$), in relation to the rotation speed ($V_r$), the greater the effect of the positive dihedral. The positive dihedral loses its effect during the autorotation phase (i.e. no forward speed $V_t=0$) (see figure \ref{fig:DiVt0}).
\begin{figure}[p]
\centering
\includegraphics[height=25em]{41_LBIEG.eps} 
\caption{Dihedral $V_t>V_r$}
\label{fig:DiVtVr}
\end{figure}

\begin{figure}[p]
\centering
\includegraphics[height=15em]{42_VT0VR.eps} 
\caption{Dihedral $V_t=0$}
\label{fig:DiVt0}
\end{figure}
\clearpage

\subsection*{The angle of attack:}
Positive angle of attack acts similar to an undercut. A positive angle means more lift and a negative angle less lift (see section \ref{fig:AngleAttack}). Increasing angle of attack increases the drag of the airfoil. 
A positive angle of attack reduces range and rotation in the final stage of flight. If the angle is too sharp, the boomerang may not quite make it back because it loses rotation during the second half of the flight. If Arm 1 has a negative angle of attack the boomerang will normally fly further. With a steep negative angle it will simply fly straight ahead. 
With a symmetrical multiblader it doesn’t matter which arm is bent. The effects just add up.

\begin{figure}[H]
\centering
\includegraphics[width=30em]{40_QBIEG.eps} 
\caption{Different angles of attack}
\label{fig:AngleAttack}
\end{figure}
\newpage

\section{Boomerang weights}
There are various techniques of weighting a boomerang. Weights or coins can be attached with tape, or embedded permanently in the boomerang. Usually, permanent weights are installed in drilled holes and filed flush with the wing. Taped on weights, in contrast to carved in weights, change the centre of rotation not only in a horizontal but also a vertical direction. The boomerang is, therefore, slightly tilted and effects similar to a warp or an under-cut area can be obtained. Special attention should be paid to how the weights are taped onto the boomerang, as there will be a change in the airfoil where the weight is attached. Because of differences in boomerangs and weighting techinques, the following effects may or may not be noticeable. 
Weights can be attached to the following points: 
\begin{center}
	\includegraphics[width=30em]{43_GEW1.eps}
\end{center}
Basic feature:
\begin{itemize}
\item further distance
\item less lay-over
\item U-shaped flight path
\end{itemize}

Variations:
\begin{itemize}
\item on top: boomerang flies higher
\item underneath: boomerang flies lower
\end{itemize}
\newpage

\begin{center}
	\includegraphics[width=30em]{44_GEW2.eps} 
\end{center}
Basic feature:
\begin{itemize}
\item further distance
\item more lay-over
\item S-shaped flight path
\end{itemize}
Variations:
\begin{itemize}
\item on top: boomerang flies lower
\item underneath: boomerang flies higher
\end{itemize}
\newpage

\begin{center}
	\includegraphics[width=30em]{45_GEWS1.eps} 
\end{center}
Basic feature: 
\begin{itemize}
\item Boomerang flies lower 
\item Boomerang keps its rotation
\item less lay-over
\end{itemize}
Variations:
\begin{itemize}
\item on top: Boomerang flies higher 
\item underneath: Boomerang flies lower, Boomerang flies further
\end{itemize}
\newpage

\begin{center}
	\includegraphics[width=30em]{46_GEWS2.eps} 
\end{center}
Basic feature:
\begin{itemize}
\item Boomerang flies higher 
\item Boomerang keeps it's rotation
\item more lay-over
\end{itemize}
Variations:
\begin{itemize}
\item on top: Boomerang flies lower, Boomerang flies further 
\item underneath: Boomerang flies higher
\end{itemize}
\newpage

\begin{center}
	\includegraphics[width=30em]{47_GEWE.eps} 
\end{center}
Basic feature:
\begin{itemize}
\item shorter distance
\end{itemize}
Variations:
\begin{itemize}
\item on top: Boomerang flies lower 
\item underneath: Boomerang flies higher
\end{itemize}
\newpage

\begin{center}
	\includegraphics[width=30em]{48_GEW12.eps} 
\end{center}
Basic feature:
\begin{itemize}
\item further distance
\end{itemize}
Variations:
\begin{itemize}
\item A,B, on top: Boomerang flies higher 
\item A,B, underneath: Boomerang flies lower 
\item A on top, B underneath: Boomerang flies	shorter distance and higher 
\item A underneath, B on top: Boomerang flies lower
\end{itemize}
\newpage

\begin{center}
	\includegraphics[width=30em]{49_GES12.eps} 
\end{center}
Basic feature:
\begin{itemize}
\item wind stabilized - unchanged distance
\end{itemize}
Variations:
\begin{itemize}
\item A, B on top: unchanged distance 
\item A, B underneath: unchanged distance 
\item A on top, B underneath: Boomerang flies a shorter distance and higher 
\item A underneath, B on top: Boomerang flies further and lower
\end{itemize}
\newpage

\begin{center}
	\includegraphics[width=30em]{50_GE12E.eps} 
\end{center}
Basic feature:
\begin{itemize}
\item Further distance 
\item Boomerang keeps its rotation 
\item Boomerang flies as though unweighted
\end{itemize}
Variations:
\begin{itemize}
\item A, B, C on top: no further change 
\item A, B, C underneath: no further change 
\item A, B on top, C underneath: Boomerang flies higher 
\item A,B, underneath, C on top: Boomerang flies further and lower 
\item A, C on top, B underneath: Boomerang flies a shorter distance and higher 
\item A on top, B, C underneath: Boomerang flies a shorter distance and higher
\item A underneath, B, C on top: Boomerang flies further and lower 
\item A, C underneath, B on top: Boomerang flies further and lower
\end{itemize}
\newpage

\section{Carving out}
To ``carve out'' means to remove substance from the topside or underside of the boomerang. Usually, either a whole arm is carved out, or just the tip is carved out. (see \ref{fig:CarveBoom}).
\begin{figure}[H]
	\centering
	\includegraphics[width=28em]{52_AUS.eps}
	\caption{Carved out boomerang profiles}
	\label{fig:CarveBoom}
\end{figure}

Carving reduces the weight of the boomerang and tends to weaken it while causing it to fly a shorter distance and higher. You can expect an extensively carved out boomerang to be more fragile and to have shorter range. If one arm is carved out significantly more than the other, the centre of rotation may shift significantly. This can affect overall lift and wind resistance.

\section{Undercut or bevel}
An undercut is an area where material has been removed from the underside of the wing from either the leading or the trailing edge. A leading edge undercut causes more lift but also increases drag whereas a trailing edge undercut reduces lift and tends to increase spin (see figure \ref{fig:LTundercut}).
\begin{figure}[H]
	\centering
	\subfigure[Trailing edge]{
	\includegraphics[width=14em]{57_BEVHI.eps}
	}
	\subfigure[Leading edge]{
	\includegraphics[width=14em]{56_BEVVO.eps}
	}
	\caption{Undercuts}
	\label{fig:LTundercut}
\end{figure}
Undercutting usually starts at the wing tip. It varies in length (see \ref{fig:undercut}). Undercuts can be used at any point on the boomerang even at the elbow.	
\begin{figure}[H]
	\centering
	\includegraphics[width=30em]{58_BEVL.eps}
	\caption{Undercut viewed from bottom}
	\label{fig:undercut}
\end{figure}
In later discussion in this book, an ``undercut'' will be referring to a 2-inch long undercut on the leading edge starting from the tip of the arm.

\section{The elbow}
	\label{sec:Elbow}
The profile of the elbow can have a large effect on the lay-over of a two-arm boomerang.

\begin{itemize}
\item An undercut on the outside edge of the elbow or a bevel on the top of the inner edge of the elbow increases the lay-over. The elbow causes lift ahead of the centre of rotation - $E_1$ positivel- and negative lift behind the centre of rotation - $E_2$ negative -. 
\item An undercut on the inner edge of the elbow or a bevel on the top of the outer edge of the elbow decrease the lay-over. The elbow causes negative lift ahead of the centre of rotation - $E_1$ negative - and a positive lift behind the centre of rotation - $E_2$ positive - (see figure \ref{fig:AirfoilElbow}).
\end{itemize}
\begin{figure}[H]
	\centering
	\includegraphics[width=30em]{42E_ELL1.eps}
	\caption{Effect of elbows wingprofile}
	\label{fig:AirfoilElbow}
\end{figure}

\section{Holes}
\label{sec:holes}
Various sizes and shapes of holes can be used anywhere on a boomerang. Generally holes increase the drag and can have other large effects on a boomerang’s flight. They can also be used for decoration or as a launching aid.

It can generally be said that:
\begin{itemize}
\item the closer the hole is to the wing tip, the slower the rotation at the end of the flight.
\item the closer the hole is to the centre of rotation, the less influence it will have on the rotation, but it does tend to reduce the speed of the boomerang.
\item The larger the hole at the centre of rotation of a multi-blader, the smaller the lay-over the boomerang will have.
\end{itemize}


Variations: (see figure \ref{fig:Holes})
\begin{itemize}
\item The position of the holes on the boomerang relative to the elbow and wingtip
\item The position of the holes on the airfoil relative to the leading and trailing edge
\item The size of the holes 
\item Sharp or rounded-off edges 
\item The slant of the holes 
\item Holes taped off on the top or underside 
\item Holes taped off on the top and underside 
\item Holes partly taped off.
\end{itemize}

\begin{figure}[H]
	\centering
	\includegraphics[width=30em]{60_LOCH.eps}
	\caption{Holes half covered}
	\label{fig:Holes}
\end{figure}

\clearpage
\section{Flaps}
A flap is a piece of tape that is stuck onto the boomerang - usually in such a way that it stands erect on the top or bottom of the boomerang. It has similar effects to a hole (see \ref{sec:holes}). The flap is much easier to attach and remove. It is easy to learn to use flaps, since they can be made and installed within a few seconds and can be easily removed and adjusted. The boomerang can then be re-tested before conditions change.

Variations: (see figure \ref{fig:Flaps})
\begin{itemize}
\item The position on the boomerang
\item The position on the profile 
\item Shape 
\item Length 
\item Direction 
\item Material
\end{itemize}

\begin{figure}[H]
	\centering
	\includegraphics[width=30em]{63_FLAPS.eps}
	\caption{Flaps mounted on wing}
	\label{fig:Flaps}
\end{figure}
\newpage

\newpage
\section{Rubberbands}
Rubber bands are easy to carry, easy to add and can largely replace major surgery on boomerang, as well as stabilize the boomerang, when thrown into the very wind. The general purpose of the rubber band is to slow the boomerang. The placement of rubber bands thus determines what is dulled. The following are just rules of thumb when using rubber bands.

\subsection{General}
The closer to the center of mass rubber bands are placed, the more it reduces the boomerang speed, and the closer rubber bands placed at the wing tips, the more brakes the rotation and reduces buoyancy.

\begin{quote}
``Rubber Bands reduces lift on the wings, helping the boomerang to fly further, because it no longer turns so sharp.'' \footnote{Jens Krabbe (Former President of Danish Boomerang Club.)}
\end{quote}

\subsection{Twobladers}
Specifically for twobladers, rubber bands put on wing 1 delay the return lay-over so that the boomerang will lie down later in it’s orbit.  Rubber bands mounted on wing 2 do the opposite, namely that the Boomerang will layover earlier in it’s course.

\subsection{Multibladers}
The multiblader’s center is in the middle and a rubber band in the center can stabilize the boomerang when throwing in the wind.  It may be useful to put rubber bands on the tips of the wings of a multiblader if it has a tendency to fly past you too fast so that you can catch the boomerang.

\newpage
\subsection{Installing rubber bands}
When mounting th rubber bands, you can do it in different ways:  You can either put them on so that they sit together, or you can wrap them around twice and spread them apart (see figure \ref{fig:Rubwing}). Spreading the rubber band creates more wind resistance resulting in a stronger effect.
You can also mount the rubber bands on tribladers on the wings and in the middle.  When you want to put a rubber band onto  the middle of a triblader, place one of the wings through the rubber band, twist it half a turn and pass the rubber band over the middle, between the two opposing wings, and then back over the first blade again (see figure \ref{fig:RubberbandTri})
\begin{figure}[H]
	\centering
	\includegraphics[width=10em]{Elastik2vinge}
	\caption{Rubberband on wing}
	\label{fig:Rubwing}
\end{figure}

\begin{figure}[p]
	\centering
		\subfigure[First step]{
			\includegraphics[height=13em]{Elastik1.eps}
		}
		\subfigure[Second step]{
			\includegraphics[height=13em]{Elastik2.eps}
		}
		\subfigure[Third step]{
			\includegraphics[height=13em]{Elastik3.eps}
		}
	\caption{Installing rubberbands on tri-blader.}
	\label{fig:RubberbandTri}
\end{figure}
\clearpage

\section{Windbreakers in general}
Windbreakers can be anything that change the behaviour of the boomerang during flight. There is no limit on ingenuity. 
Experiment on the following windbreakers:
\begin{itemize}
\item Rubberbands 
\item Bandages
\item A complete roll of cellotape
\item String 
\item Wire
\item $\cdots$
\end{itemize}

\section{Slits}
Slits are, in principle, nothing more than elongated holes. They can be narrow or wide, long or short, profiled or unprofiled in various positions on the boomerang. Adam ``Schlitzer'' Müller developed the slotting and comb (see next section) technology used in competition boomerangs today. 

The wing is, in effect, divided into two small wings with a wide profiled slit between them. The boomerang tends to be more stable in wind and retains most of its spin but loses much of its forward speed. Be careful! Avoid weakening the wings excessively. Polypropylene, fibreglass and carbon fibre are generally the best materials for making slotted boomerangs.
\newpage

\section{Combs}
Combs are thin slits on the front or back of the airfoil (see figure \ref{fig:Combs}). They cause slight turbulence, similar to flaps or holes but the effect is more even. 

The boomerang flies further with a comb on the trailing edge without losing much rotation. It is not necessarily more stable in wind. If the comb is on the leading edge, stability in wind is improved. The closer the comb is to the tip, the more the rotation is reduced. The thicker the airfoil, the greater the effect.

The length and width of the comb area that is practical, depends on the strength of the material. But be warned! Because of the combs, the wing loses strength and is much more easily broken. Polypropylene, Fibre glass and Carbon fibre are the best materials.

\begin{figure}[H]
	\centering
	\includegraphics[width=30em]{66_KAMM.eps}
	\caption{Combs cut into the airfoil}
	\label{fig:Combs}
\end{figure}
